\SourceLeaf{334}%
\[
\begin{aligned}
\log t={}&\log\tau+\frac{Mk^4}{12aa(1-ee)}\mathord{.}rr+\tfrac1{24}M\tau\tau\\
&+\frac M{24k^4}(5ee+(4ee-14e^4)ss+5e^4s^4)\beta\beta\\
\log b={}&\log\beta+\frac{Mk^4}{12aa(1-ee)}\mathord{.}rr+\tfrac18M\tau\tau\\
&-\frac M{24k^4}(2+ee-(8ee-14e^4)ss-9e^4s^4)\beta\beta\\
\log l={}&\log\lambda+\tfrac1{24}M\tau\tau-\frac{(1-ee)M}{24k^4}(1-10eess)\beta\beta
\end{aligned}
\]

Da, wie man leicht sieht, in allen bisher entwickelten Formeln die Grössen \(t,\tau,b,\beta,l,\lambda\) als in Theilen des Halbmessers ausgedrückt angenommen sind, so wird man, wenn jene in Secunden ausgedrückt und dieselben Bezeichnungen für sie beibehalten werden sollen, den Formeln für \(\tau,\beta,\lambda\) (Art. 29) noch den Factor \(\frac1\rho\) beifügen müssen; in den Gleichungen 9, 10, 11 hingegen, so wie in den daraus abgeleiteten logarithmischen, muss den Gliedern, die \(\tau\tau\) oder \(\beta\beta\) enthalten, noch der Factor \(\rho\rho\) zugesetzt werden, wo \(\rho\) (eben so wie oben Art. 16 und 19) die Grösse des Bogens von einer Secunde in Theilen des Halbmessers bedeutet. Behält man nun auch noch \(\mu\) in der oben gebrauchten Bedeutung bei, nemlich
\[
\mu=\tfrac1{12}M\rho\rho
\]
und schreibt zur Abkürzung
\[
\begin{aligned}
(1)&=\frac{k}{a\rho}\\
(2)&=\frac{k^3}{a(1-ee)\rho}\\
(3)&=\frac{Mk^4}{12aa(1-ee)}\\
(4)&=\frac{\mu}{2k^4}(5ee+(4ee-14e^4)ss+5e^4s^4)\\
(5)&=\frac{\mu}{2k^4}(2+ee-(8ee-14e^4)ss-9e^4s^4)\\
(6)&=\frac{(1-ee)\mu}{2k^4}(1-10eess)\\
(7)&=\tfrac12\mu
\end{aligned}
\]
so ist unsre Auflösung in folgenden sechs Formeln enthalten:
\[
\begin{aligned}
\tau&=(1)r\sin T\operatorname{tang}B\\
\beta&=(2)r\cos T\\
\lambda&=(1)r\sin T\sec B\\
\log t&=\log\tau+(3)rr+(4)\beta\beta+(7)\tau\tau\\
\log b&=\log\beta+(3)rr-(5)\beta\beta+3(7)\tau\tau\\
\log l&=\log\lambda-(6)\beta\beta+(7)\tau\tau
\end{aligned}
\]
